\chapter[Hindi]{\textsanskrit{हिंदी}} \label{chap:hi}
{\color{intro}\bt{\textsanskrit{जो सबको अपनाकर भी अलग रही, जिसमें परमात्मा का प्राचीन तत्त्व और कलियुग का मैल दोनों वास करते हैं, और जिसकी भिन्नता उसे स्वीकारने वाले विशाल भारत गणराज्य का सामाजिक प्रतीक है, ऐसी योगिसदृश भाषा की वास्तविक विशेषता उसके भीतर स्थित ज्ञान में ही पाई जा सकेगी।}}}

\textsanskrit{हिंदी} is a language which attempts to best represent the Indian spirit, yet fails to depict its true glory due to its ridiculous range. This is a country whose states are often larger than entire European nations, has different languages and  cultures per state, and boasts of a population which is so diverse, that it consists of members belonging to the deepest pit of ignorance till the highest pinnacle of virtue. Were the diversity of the entire world be crystallized and compared against this humoungous subcontinent known as Bharat (\textsanskrit{भारत}), I dare say that \textsanskrit{भारत} will win, and it is not a close competition.

Practically speaking, \textsanskrit{हिंदी} behaves somewhat like English: its origins, like with most Indian languages can be traced back to \textsanskrit{संस्कृतम्}, but due to historical and official reasons, it has attempted, similar to English to integrate as much diversity within itself as possible, giving rise to similar problems of purity and consistency that English faces, but in a different and arguably more useful form. As a result, \textsanskrit{हिंदी} is to \textsanskrit{भारत} what English is to the world: a generally accepted universal language. 

Learning \textsanskrit{हिंदी} is like learning multiple languages that disguise themselves as one, but provides the advantage of a shared vocabulary, along with economic access to the  huge Indian market.

An important point must be addressed about purity before proceeding. Due to its course of development, pure \textsanskrit{हिंदी} doesn't really exist and is rather impractical. It is thus a reasonable compromise in this case  to accept foreign words that at least sound like they fit in with everything else as `pure'.

The \textsanskrit{देवनागरी} (Devanagari) script that \textsanskrit{हिंदी} uses  has NO exceptions in pronunciation. Alchemically, the script takes some getting used to as an equivalent exchange. An additional  advantage that \textsanskrit{हिंदी} offers is spelling flexibility: Because of extreme  lingual intermixing, words often do not have unique representations. Any spellings that sound similar are treated as equivalents:
\begin{example}
    \textsanskrit{सम्भव / संभव ; हिन्दी / हिंदी ; पडना / पड़ना ; फर्क / फरक / फ़र्क़}
\end{example}
This mostly happens with  \textsanskrit{न्} and  \textsanskrit{म्}, which can be converted to an \textsanskrit{अनुस्वार} (\textsanskrit{ं}).

No comment on the roughness of this perspective is appropriate, since it is as rough in \autoref{sec:histce} as it is  smooth in \autoref{sec:hinoun}. This interesting mix makes it a good potential medium for solving other perspectives, but those considerations are out of the scope of this English work.

\section[Noun]{Noun (\textsanskrit{संज्ञा})} \label{sec:hinoun}
\section[Verb]{Verb (\textsanskrit{क्रिया})} \label{sec:hiverb}

\section[Preposition]{Preposition (\textsanskrit{सम्बन्धबोधक})} \label{sec:hiprep}
\section[Sentence]{Sentence (\textsanskrit{वाक्य})} \label{sec:histce}
\section[Specific constructions]{Specific constructions (\textsanskrit{अनोखी रचनाएँ})} \label{sec:hispco}
\section[Closing remarks]{Closing remarks (\textsanskrit{अंतिम टिप्पणियाँ})} \label{sec:hirmrk}